% Section 12.2, from lectures/L12/appendix/b_bit_timer.md.

\section{Designing \code{bit_timer}}
\label{c12:sec:bittimer}

\code{bit_timer} divides the 50 MHz system clock into CAN bit periods and produces two one-cycle
pulses that the rest of the controller runs on: one marking where in the bit the bus is to be
sampled, one marking the bit's end.

Why a bit period is divided up at all, and why the sample point sits late in it rather than in the
middle, is in \canref{5}; this section builds the module that does what that chapter describes, with
a fixed sample point instead of configurable segments.

\subsection{Interface}
\label{c12:sub:btiface}

\modulefig[0.62]{lectures/L12/appendix/images/bit_timer.png}

\begin{booktable}{@{}p{5.6em}p{3.4em}p{5.2em}L@{}}
  \thead{Port} & \thead{Dir.} & \thead{Type} & \thead{Meaning} \\
  \midrule
  \code{clock} & in & \code{std_logic} & 50 MHz system clock. \\
  \code{reset_s2_n} & in & \code{std_logic} & Active-low reset, synchronised with two flip-flops. \\
  \code{enable} & in & \code{std_logic} & Runs the counter while \code{'1'}; holds it while
    \code{'0'}. \\
  \code{resync} & in & \code{std_logic} & Forces the counter back to tick 0 at the next clock
    edge. \\
  \code{sample} & out & \code{std_logic} & Flag indicating that the bus is to be sampled. \\
  \code{bit_done} & out & \code{std_logic} & Flag indicating that this bit is over and the next one
    begins. \\
\end{booktable}

No generics. \code{bit_timer_tb} binds its ports \textbf{positionally}, so the order and the types
above have to match exactly. The names are yours; keeping these makes the book's text and your code
speak about the same signals.

From \code{can_def} (\chapref{c10:ch}) it needs \code{TICKS_PER_BIT} (= 50, the ticks in a bit
period) and \code{SAMPLE_TICK} (= 35, that is 70~\% in, late enough for the bus to have settled).

\subsection{Behaviour}
\label{c12:sub:btbehav}

A single free-running counter over \code{0 .. TICKS_PER_BIT - 1} is enough. At every clock edge:
default both pulses to \code{'0'} and then, in priority order:
\begin{enumerate}
  \item \textbf{\code{reset_s2_n = '0'}}: clear the counter and both pulses.
  \item \textbf{\code{resync = '1'}}: restart the bit period; put the counter back to \code{0}. This
    wins over \code{enable}, and is used once per frame, at SOF, so that this node's bit timing
    starts afresh, aligned with the edge opening the frame rather than wherever an earlier count
    happened to stand.
  \item \textbf{\code{enable = '1'}}: pulse \code{sample} during the single tick where
    \code{counter = SAMPLE_TICK}. At \code{counter = TICKS_PER_BIT - 1}, set \code{bit_done} and let
    the counter wrap to \code{0}; otherwise increment it.
  \item \textbf{Otherwise (disabled)}: hold the counter, both pulses low.
\end{enumerate}

This is the same counter-with-priority-branches idiom as \chapref{c7:ch}'s \code{timer}, only with
two pulse outputs instead of one.

\begin{aside}
A note about \code{enable}, since it is the one port whose eventual use is not obvious. The natural
guess is that \code{can_controller} holds the timer while the bus is idle, and it turns out not to:
\chapref{c16:ch}'s bus-idle rule counts recessive \textbf{bit periods}, so the timer has to keep
running in \code{STATE_IDLE} too, and the finished design ties \code{enable} high in every state.
Build the port anyway; the testbench checks it, and a timer that cannot be stopped is a worse building
block than one that can.
\end{aside}

\subsection{One bit period, tick by tick}
\label{c12:sub:bttrace}

The rules above are easier to trust once seen on a time axis. With \code{TICKS_PER_BIT = 50} and
\code{SAMPLE_TICK = 35}, an enabled bit period looks like this. Every column is a rising clock edge;
\code{counter} is the value the process sees at that edge, and the two rows below it are the pulses
that edge produces.

\begin{textdrawing}
counter      0   1   2  ..  34  35  36  ..  48  49   |   0   1  ..
sample       0   0   0  ..   0   1   0  ..   0   0   |   0   0  ..
bit_done     0   0   0  ..   0   0   0  ..   0   1   |   0   0  ..
\end{textdrawing}

Three things to read out of that. Both outputs are \textbf{one tick wide}, not levels. \code{sample}
comes at tick 35 and \code{bit_done} at tick 49, so \textbf{\code{sample} always fires first, 14
ticks earlier in the same period}; a receiver reads the bus 70~\% in and only after that does the
period end. And the period boundary is the wrap from 49 back to 0, which is the same edge
\code{bit_done} marks.

\textbf{Now the same picture with a \code{resync}.} Suppose the timer is 20 ticks into a period when
SOF arrives and \code{can_controller} pulses \code{resync}:

\begin{textdrawing}
tick        19  20  21  22  ..  55  56  ..  69  70  71
resync       0   1   0   0  ..   0   0  ..   0   0   0
counter     19  20   0   1  ..  34  35  ..  48  49   0
sample       0   0   0   0  ..   0   1  ..   0   0   0
bit_done     0   0   0   0  ..   0   0  ..   0   1   0
\end{textdrawing}

The counter is thrown away mid-count and starts over, so \code{counter} and \code{tick} no longer
agree. What matters is the sequence: \code{bit_done} arrives at tick \textbf{70}, a whole
\code{TICKS_PER_BIT} after the resynchronisation, rather than at tick 49 where the interrupted period
would have ended. The node has adopted the transmitter's bit boundary instead of its own, which is
the whole point of \code{resync} and exactly what the testbench's fourth case measures.

Note as well that \code{resync} restarts the period even though \code{enable} is high throughout;
that is why it sits above \code{enable} in the priority order.

\subsection{What the testbench pins down}
\label{c12:sub:bttb}

\begin{itemize}
  \item Held in reset, \code{sample} and \code{bit_done} stay at \code{'0'}.
  \item While \code{enable = '0'} the counter never advances; both pulses stay at \code{'0'}.
  \item While enabled, \code{sample} pulses exactly at \code{SAMPLE_TICK} and \code{bit_done} exactly
    at the last tick, once per period.
  \item After a \code{resync} mid-period, the next \code{bit_done} comes exactly
    \code{TICKS_PER_BIT} ticks later, not where the interrupted period would have ended.
\end{itemize}

\subsection{Putting it into \code{can_controller}}
\label{c12:sub:btwire}

\code{bit_timer} joins the two \code{meta_prev} instances from \secref{c12:sec:metaprev} inside the
top level written in \chapref{c11:ch}, and is the first block in there that knows anything about CAN.
Instantiating it is the same two-part change to \file{controller/can_controller.vhd} as
\code{meta_prev} was, and it takes \code{reset_s2_n} from the synchroniser you just wired in rather
than the port \code{reset_n}.

First declare the four signals the instance connects to, in the architecture's declarative part:

\begin{vhdlcode}
    -- bit_timer.
    signal bt_enable, bt_resync, bt_sample, bt_bit_done : std_logic;
\end{vhdlcode}

Then instantiate the entity itself, after \code{begin}:

\begin{vhdlcode}
    bit_timer1: entity work.bit_timer
        port map(clock, reset_s2_n, bt_enable, bt_resync, bt_sample, bt_bit_done);
\end{vhdlcode}

Nothing drives \code{bt_enable} or \code{bt_resync} yet, and nothing reads \code{bt_sample} or
\code{bt_bit_done}; the state machine doing both is \chapref{c16:ch}'s. That is expected at this
point. What matters now is that \file{can_controller.vhd} still analyses cleanly, which \code{make
build} checks as soon as the file exists, and that \code{bit_timer} itself passes
\code{bit_timer_tb}.
